\documentclass[10pt,a4paper]{article}
\usepackage[utf8]{inputenc}
\usepackage[T1]{fontenc}
\usepackage[margin=1.8cm,top=2.6cm,bottom=2.2cm]{geometry}
\usepackage{parskip}
\usepackage{microtype}
\usepackage{amsmath,amssymb,mathtools}
\usepackage{array,booktabs,tabularx,multirow}
\usepackage[table]{xcolor}
\usepackage{graphicx,float}
\usepackage{titlesec}
\usepackage{enumitem}
\usepackage{fancyhdr}

% ---- Section numbering: Roman numerals ---------------------------------
\renewcommand\thesection{\Roman{section}}
% ---- Equation numbering: (I.1), (II.2), etc. ---------------------------
\numberwithin{equation}{section}

% ---- Section formatting ------------------------------------------------
\titleformat{\section}{\large\bfseries}{\thesection.\;}{0em}{}[\titlerule]
\titleformat{\subsection}{\normalsize\bfseries}{}{0em}{}
\titleformat{\subsubsection}{\normalsize\itshape}{}{0em}{}
\titlespacing{\section}{0pt}{1.6ex plus .2ex}{0.6ex}
\titlespacing{\subsection}{0pt}{1.0ex}{0.3ex}

% ---- Page headers ------------------------------------------------------
\pagestyle{fancy}
\fancyhf{}
\setlength{\headheight}{15pt}
\fancyhead[L]{\small\textit{Layer~9: Surface Enzyme Kinetics}}
\fancyhead[R]{\small Aaran Nihalani $\bullet$ Eton iGEM Dry Lab}
\fancyfoot[C]{\small\thepage}
\renewcommand{\headrulewidth}{0.35pt}

% ---- Macros ------------------------------------------------------------
\newcommand{\Mms}{\text{M}^{-1}\text{s}^{-1}}
\newcommand{\cit}[1]{[#1]}
\newcommand{\kcat}{k_{\text{cat}}}
\newcommand{\keff}{k_{\text{eff}}}
\newcommand{\GS}{\Gamma_S}
\newcommand{\GES}{\Gamma_{ES}}
\newcommand{\GP}{\Gamma_P}
\newcommand{\Gstar}{\Gamma^{*}}
\newcommand{\Gz}{\Gamma_0}
\newcommand{\Da}{\mathrm{Da}}

\begin{document}

%=======================================================================
\begin{center}
{\LARGE\bfseries Layer 9: Surface Enzyme Kinetics}\\[0.4em]
{\large Cleavage Kinetics of Surface-Immobilised SGKIPRTLTA by MMP-9}\\[0.25em]
{\small\color{gray} Eton iGEM 2026, Dry Lab \quad$|$\quad
Layer~9 of 17-Layer Biosensor Pipeline \quad$|$\quad
\texttt{surface\_kinetics.py}}
\end{center}
\vspace{0.6em}

%=======================================================================
\section{Physical Picture}
\label{sec:physical}

The biosensor probe SGKIPRTLTA is surface-immobilised on the gold electrode via a
Cys-thiol anchor and carries a methylene blue (MB) redox label at the C-terminal
lysine. At surface concentration $\Gz$ (mol~cm$^{-2}$), the intact probe keeps MB
tethered close to the electrode; MMP-9 cleavage releases the MB-bearing fragment
into solution, decreasing the SWV peak current. The fractional cleavage $F_k(t)$
is therefore the central output of this layer and the direct precursor of the
ratiometric signal $\Delta R/R_0$ (Layer~13).

The governing kinetics are heterogeneous: the enzyme ($E$, MMP-9) diffuses freely
in solution while the substrate ($S$, probe) is fixed at the surface. Classical
solution-phase Michaelis--Menten equations, derived under the condition
$[S] \gg [E]$, do not apply in this configuration because the total molar quantity
of surface-immobilised peptide ($\Gz \cdot A \sim 10$~fmol) can be comparable to
or less than the bulk enzyme content. Anne and Demaille~\cit{1} derived exact
kinetic equations for this heterogeneous surface mechanism without invoking the
steady-state approximation.

The enzyme input $[E]_t = [\text{MMP-9}]_{\text{free}}$ is taken directly from the
Layer~0 wound biology ODE output: 30~nM (chronic wound, 30-min equilibration) and
6~nM (healing wound). TIMP-1 inhibition is fully accounted for in these values;
no further correction is needed here.

%=======================================================================
\section{Three-State Surface ODE System}
\label{sec:ode}

\subsection{Reaction mechanism}

The surface enzyme reaction follows the Michaelis--Menten mechanism with the
substrate anchored to the surface~\cit{1}:
\begin{equation}
  S + E \;\underset{k_{-1}}{\overset{k_1}{\rightleftharpoons}}\; ES
           \;\xrightarrow{\;\kcat\;}\; P + E
  \label{eq:mechanism}
\end{equation}
where $S$, $ES$, $P$ are surface species with concentrations $\GS$, $\GES$,
$\GP$ (mol~cm$^{-2}$); $k_1$ (M$^{-1}$s$^{-1}$) is the second-order binding rate
constant; $k_{-1}$ (s$^{-1}$) is the dissociation rate constant; and $\kcat$
(s$^{-1}$) is the catalytic rate constant. The enzyme concentration at the
electrode surface $[E]_t$ (M) is treated as constant throughout (justified in
Section~\ref{sec:dimensionless}).

\subsection{Governing ODEs}

The time evolution of each surface species is~\cit{1} (Eq.~1):
\begin{equation}
  \label{eq:odes}
  \frac{d\GS}{dt}  = k_{-1}\,\GES - k_1[E]_t\,\GS, \qquad
  \frac{d\GES}{dt} = k_1[E]_t\,\GS - (k_{-1}+\kcat)\,\GES, \qquad
  \frac{d\GP}{dt}  = \kcat\,\GES
\end{equation}
with surface conservation and initial conditions:
\begin{equation}
  \label{eq:conservation}
  \GS + \GES + \GP = \Gz, \qquad
  \GS(0) = \Gz,\quad \GES(0) = \GP(0) = 0
\end{equation}
Introducing the unconverted substrate $\Gstar = \GS + \GES$ and defining
$k_1^* = k_1[E]_t$, the system reduces to the two-variable form used by
Anne and Demaille~\cit{1} for the Laplace analysis:
\begin{equation}
  \label{eq:reduced}
  \frac{d\GES}{dt} = k_1^*\,\Gstar - (k_1^* + k_{-1} + \kcat)\,\GES,
  \qquad
  \frac{d\Gstar}{dt} = -\kcat\,\GES
\end{equation}

%=======================================================================
\section{Analytic Solutions}
\label{sec:analytic}

\subsection{Steady-state approximation}

Setting $d\GES/dt = 0$ in Eq.~\eqref{eq:reduced} and substituting the conservation
constraint yields the steady-state progress curve~\cit{1} (Eq.~3):
\begin{equation}
  \label{eq:ss}
  \theta_P^{\text{ss}} = \frac{\GP}{\Gz} = 1 - \exp(-\keff\,t), \qquad
  \keff = \frac{\kcat}{1 + K_M/[E]_t}, \quad
  K_M = \frac{k_{-1} + \kcat}{k_1}
\end{equation}
This is the Langmuirian (LG) progress curve. In the linear regime
($[E]_t \ll K_M$): $\keff \approx (\kcat/K_M)[E]_t \equiv \keff^{\text{surf}}[E]_t$,
giving a simple pseudo-first-order decay.

\subsection{Exact biexponential solution}

The reduced system Eq.~\eqref{eq:reduced} is linear and solved exactly via Laplace
transform~\cit{1}. Define the composite parameters:
\begin{equation}
  \label{eq:composite}
  k_1^* = k_1[E]_t, \qquad
  \sigma = k_1^* + k_{-1} + \kcat, \qquad
  \Delta = \sigma^2 - 4k_1^*\kcat > 0
\end{equation}
The exact solutions are~\cit{1} (Eqs.~4--5):
\begin{align}
  \GES(t) &= \frac{k_1^*\,\Gz}{\sqrt{\Delta}}
    \left[\exp\!\left(-\frac{\sigma-\sqrt{\Delta}}{2}\,t\right)
         -\exp\!\left(-\frac{\sigma+\sqrt{\Delta}}{2}\,t\right)\right]
  \label{eq:Ges} \\[0.5em]
  \Gstar(t) &= \frac{\Gz}{2\sqrt{\Delta}}
    \left[(\sigma+\sqrt{\Delta})\exp\!\left(-\frac{\sigma-\sqrt{\Delta}}{2}\,t\right)
         -(\sigma-\sqrt{\Delta})\exp\!\left(-\frac{\sigma+\sqrt{\Delta}}{2}\,t\right)\right]
  \label{eq:Gstar}
\end{align}
The fractional product coverage (fractional cleavage) is then:
\begin{equation}
  \label{eq:Fk_biexp}
  F_k(t) = \theta_P(t) = 1 - \frac{\Gstar(t)}{\Gz}
\end{equation}

\subsection{$\lambda$ parameter and dimensionless universal form}

Anne and Demaille~\cit{1} introduce the dimensionless kinetic zone parameter:
\begin{equation}
  \label{eq:lambda}
  \lambda = \frac{\sigma}{\sqrt{\Delta}}
  = \frac{1}{\sqrt{1 - 4k_1^*\kcat/\sigma^2}}
  = \frac{1}{\sqrt{1 - 4k_1\kcat[E]_t\,/\,(k_1^*+k_{-1}+\kcat)^2}}
  \geq 1
\end{equation}
and the dimensionless time $\tau = \sigma(1 - 1/\lambda)\,t$. In these variables the
exact progress curve Eq.~\eqref{eq:Fk_biexp} becomes~\cit{1} (Eq.~7):
\begin{equation}
  \label{eq:theta_exact}
  \theta_P(\tau,\lambda) = 1 - \frac{1}{2}\!\left\{
    (1+\lambda)\exp\!\left[-\frac{\tau}{2}\right]
    +(1-\lambda)\exp\!\left[-\frac{1}{2}\frac{\lambda+1}{\lambda-1}\,\tau\right]
  \right\}
\end{equation}
This single equation captures the entire range of possible kinetic behaviours
through the value of $\lambda$ alone~\cit{1}. Two limiting cases are:
\begin{itemize}[noitemsep]
  \item \textbf{LG limit} ($\lambda \to 1$): Eq.~\eqref{eq:theta_exact}
    reduces to $\theta_P = 1 - \exp(-\tau/2)$, a simple monoexponential.
    The steady-state approximation is exact.
    Diagnostic: $t_{1/2}/t_{1/4} \to 2.41$~\cit{1}.
  \item \textbf{FS limit} ($\lambda \gg 1$): Eq.~\eqref{eq:theta_exact}
    reduces to~\cit{1} (Eq.~8):
    \begin{equation}
      \label{eq:fs}
      \theta_P = 1 - \!\left(1 + \frac{\tau}{2}\right)\exp\!\left(-\frac{\tau}{2}\right)
    \end{equation}
    The progress curve is S-shaped (sigmoidal foot at short times).
    The steady-state approximation overestimates $\theta_P$ and fails here.
    Diagnostic: $t_{1/2}/t_{1/4} \to 1.75$~\cit{1}.
\end{itemize}
The zone diagram~\cit{1} (Fig.~4) maps these regimes in
$\log(k_{-1}/\kcat) - \log([E]_t/K_M)$ space. The LG zone applies when
$k_{-1}/\kcat \gg \Phi$, where $\Phi = 4[E]_t/K_M\,/(1+[E]_t/K_M)^2$, or
unconditionally when $[E]_t/K_M \leq 0.2$ or $[E]_t/K_M \geq 5$.

%=======================================================================
\section{Surface Probe Saturation: The Integrated MM Equation}
\label{sec:surface_sat}

\subsection{A second, independent saturation condition}

The A\&D analysis in Sections~\ref{sec:ode}--\ref{sec:analytic} characterises
the kinetic zone along the \emph{enzyme} axis: does the enzyme--probe complex
form and dissociate fast enough that the steady-state approximation holds?
The conclusion $\lambda \approx 1$ (Section~\ref{sec:regime}) is a statement
about enzyme-side kinetics and holds independently of the probe surface coverage.

A second condition must be checked separately: whether the probe density
$\Gz$ is large or small relative to the \emph{surface} Michaelis constant
$K_{m,\text{surf}}^\Gamma$ (units: mol~cm$^{-2}$). The A\&D Eqs.~4--7 assume
the probe is always in the unsaturated (dilute) regime, i.e.\ effectively
$K_{m,\text{surf}}^\Gamma \to \infty$. At the surface coverages used in our
device this assumption requires verification.

\subsection{The heterogeneous surface MM ODE with correct units}

The surface Michaelis--Menten equation governing probe depletion, with
substrate and product both expressed as surface concentrations
(mol~cm$^{-2}$) and enzyme as a volume concentration (mol~L$^{-1}$), is
(Squires et~al.\ 2008~\cit{4}, Eq.~1):
\begin{equation}
  \frac{d\GS}{dt} = -\frac{\kcat\,[E]_t\,\GS}{K_{m,\text{surf}}^\Gamma + \GS}
  \label{eq:mm_surface}
\end{equation}
In the A\&D notation, $K_{m,\text{surf}}^\Gamma$ is the probe density at
which the per-probe turnover rate is half its maximum value. Its units are
mol~cm$^{-2}$, not mol~L$^{-1}$.

\subsection{Conversion from solution $K_m$ to surface units}

The solution-phase $K_m^{\text{sol}}$ (mol~L$^{-1}$, confirmed for
SGKIPRTLTA from Kridel~2001~\cit{2}) is related to the surface saturation
constant by the effective thickness of the accessible probe layer $\delta_{\text{eff}}$
(Layer~8):
\begin{equation}
  K_{m,\text{surf}}^\Gamma
  = K_m^{\text{sol}}\,[\text{mol~L}^{-1}]
    \times 10^{-3}\,[\text{L~cm}^{-3}]
    \times \delta_{\text{eff}}\,[\text{cm}]
  \label{eq:Kmsurf}
\end{equation}
For $K_m^{\text{sol}} = 1.28$~mM and $\delta_{\text{eff}} = 5$~nm $= 5\times10^{-7}$~cm
(Layer~8 probe extent, unconfirmed):
\begin{equation}
  K_{m,\text{surf}}^\Gamma
  = 1.28\times10^{-3} \times 10^{-3} \times 5\times10^{-7}
  = 6.4\times10^{-13}\;\text{mol~cm}^{-2}
  \label{eq:Kmsurf_val}
\end{equation}

\subsection{The probe-saturation ratio $\rho$}

Define the dimensionless probe-saturation parameter:
\begin{equation}
  \rho = \frac{\Gz}{K_{m,\text{surf}}^\Gamma}
  \label{eq:rho}
\end{equation}
$\rho$ controls the shape of the progress curve independently of $\lambda$
(which controls the enzyme-saturation regime). Two limits:
\begin{itemize}[noitemsep]
  \item $\rho \ll 1$: probe is always in the unsaturated regime
    $\GS \ll K_{m,\text{surf}}^\Gamma$; Eq.~\eqref{eq:mm_surface} reduces to
    first-order in $\GS$; the monoexponential Eq.~\eqref{eq:monoexp} is
    recovered with $\keff^{\text{surf}} = \kcat/K_m^{\text{sol}}$.
  \item $\rho \gg 1$: the probe is initially in the zeroth-order (saturated)
    regime; the cleavage rate is independent of $\GS$ and proportional only
    to $[E]_t$: $d\GS/dt = -\kcat[E]_t$ → $F_k = \kcat[E]_t\,T/\Gz$.
    The progress curve is \emph{linear} in time, with a Langmuir-like saturation
    shape at late times as $\GS$ falls through $K_{m,\text{surf}}^\Gamma$.
  \item $\rho \sim 1$: intermediate (physically realistic) case: the progress
    curve is \emph{between} the monoexponential and the Langmuir-like linear
    regime, as treated by the integrated equation below.
\end{itemize}

At the unconfirmed reference values ($\Gz = 2\times10^{-12}$~mol~cm$^{-2}$
from Pheeney~2012~\cit{5} midpoint; no steric correction):
$\rho_0 = 2\times10^{-12}/6.4\times10^{-13} \approx 3.1$.
With the Layer~8 steric correction (which increases $K_{m,\text{surf}}^\Gamma$
by $f_{\text{steric}}^{-1}$; unconfirmed factor $\approx 2$--$10\times$):

\begin{table}[H]
\centering
\small\renewcommand{\arraystretch}{1.25}
\begin{tabular}{@{}cccc@{}}
\toprule
$f_{\text{steric}}$ (Layer~8) &
$K_{m,\text{surf}}^\Gamma$ (mol~cm$^{-2}$) &
$\rho$ &
Dominant regime \\
\midrule
0.1 (high crowding) & $6.4\times10^{-12}$ & 0.31 & near first-order \\
0.3                 & $2.1\times10^{-12}$ & 0.95 & intermediate \\
0.5                 & $1.3\times10^{-12}$ & 1.5  & intermediate \\
1.0 (no correction) & $6.4\times10^{-13}$ & 3.1  & intermediate--zeroth \\
\bottomrule
\end{tabular}
\caption{Probe-saturation ratio $\rho$ across the plausible steric
correction range from Layer~8. $\rho = 1$ is the critical point at which
the progress curve deviates most from the monoexponential.
The system is unlikely to be in the pure monoexponential limit ($\rho \ll 1$)
at clinical probe densities unless crowding is severe ($f_{\text{steric}} < 0.05$).}
\label{tab:rho}
\end{table}

\subsection{Implicit integrated Michaelis--Menten solution}

Integrating Eq.~\eqref{eq:mm_surface} from $\GS = \Gz$ to
$\GS = \Gz(1-F_k)$, with $[E]_t$ constant (justified by $\Da \ll 1$,
Section~\ref{sec:dimensionless}):
\begin{equation}
  \boxed{
  \kcat\,[E]_t\,t
  = K_{m,\text{surf}}^\Gamma\,\ln\!\frac{1}{1-F_k}
  + \Gz\,F_k
  }
  \label{eq:imp_mm}
\end{equation}
This is the primary model for $F_k(T)$. It is an implicit equation in $F_k$
and must be inverted numerically (bisection or Newton's method, implemented in
\texttt{surface\_kinetics.py}). In dimensionless form, with
$\tilde{t} = \kcat[E]_t\,t/K_{m,\text{surf}}^\Gamma$ and $u = F_k$:
\begin{equation}
  \tilde{t} = -\ln(1-u) + \rho\,u
  \label{eq:imp_mm_dimless}
\end{equation}
Recovery of the two limits:
\begin{align}
  \rho = 0:\quad &\tilde{t} = -\ln(1-u)
    \;\Rightarrow\; F_k = 1 - \exp\!\left(-\frac{\kcat}{K_m^{\text{sol}}}[E]_t\,t\right)
    \quad \text{(monoexponential)} \label{eq:limit_mono}\\[0.3em]
  \rho \to \infty:\quad &\tilde{t} \approx \rho\,u
    \;\Rightarrow\; F_k = \frac{\kcat[E]_t\,t}{\Gz}
    \quad \text{(Langmuir-like, linear in }t\text{)} \label{eq:limit_lang}
\end{align}
The two-parameter form $\{\kcat[E]_t,\,K_{m,\text{surf}}^\Gamma\}$ (equivalently
$\{k_{\text{obs}},\,\rho\}$) is fit to experimental progress curves from a
single incubation time-series. The monoexponential (one-parameter fit) is
recovered as the $\rho \to 0$ special case.

\subsection{Short-time approximation and discrimination ratio}

At early times ($F_k \ll 1$), Taylor-expanding Eq.~\eqref{eq:imp_mm_dimless}:
$\tilde{t} \approx (1+\rho)\,u + \ldots$, so:
\begin{equation}
  F_k(T) \approx \frac{\kcat[E]_t\,T}{K_{m,\text{surf}}^\Gamma + \Gz}
  = \frac{\kcat[E]_t\,T}{K_{m,\text{surf}}^\Gamma(1 + \rho)}
  \quad \text{(short-}T\text{ linear regime)}
  \label{eq:short_time}
\end{equation}
The ratio of chronic to healing signal at short times is independent of both
$\kcat$ and $K_{m,\text{surf}}^\Gamma$ (and of $\rho$):
\begin{equation}
  \frac{F_k^{\text{chronic}}}{F_k^{\text{healing}}} \approx
  \frac{[E]_{\text{chronic}}}{[E]_{\text{healing}}} = 5\times
  \label{eq:ratio_rho}
\end{equation}
This result, previously derived from the monoexponential limit, holds
for \emph{all} $\rho$ in the short-time regime. The 5$\times$ discrimination
ratio is therefore robust to the value of $\rho$ at early incubation times.

%=======================================================================
\section{Kinetic Regime at Clinical MMP-9 Concentrations}
\label{sec:regime}

\subsection{$\lambda \approx 1$: the Langmuirian limit is guaranteed}

At clinical concentrations $[E]_t \in [0.19,\,2.17]$~nM and any physically
realisable binding rate constant $k_1 \leq 10^8\ \Mms$ (diffusion limit for
a 92~kDa enzyme):
\begin{equation}
  \label{eq:k1star}
  k_1^* = k_1[E]_t \leq 10^8 \times 2.17\times10^{-9} = 0.22\ \text{s}^{-1}
\end{equation}
The solution-phase value $\kcat^{\text{sol}} = 86.4$~s$^{-1}$~\cit{2} is an upper
bound for the surface $\kcat$. In the regime $k_1^* \ll \kcat$:
\begin{equation}
  \label{eq:lambda_approx}
  \Delta \approx (k_{-1}+\kcat)^2 + 2k_1^*(k_{-1}-\kcat) \approx (k_{-1}+\kcat)^2,
  \quad
  \lambda = \frac{k_1^*+k_{-1}+\kcat}{\sqrt{\Delta}} \approx 1 + O\!\left(\frac{k_1^*}{\kcat}\right)
\end{equation}
Since $k_1^*/\kcat^{\text{sol}} \leq 2.5\times10^{-3}$, we have $\lambda - 1 < 3\times10^{-3}$
regardless of the unknown surface $k_{-1}$. The system is firmly in the LG zone.

\subsection{Collapse to monoexponential}

In the LG limit, the exact biexponential Eq.~\eqref{eq:theta_exact} degenerates
to the Langmuirian monoexponential irrespective of the $k_{-1}/\kcat$ ratio:
\begin{equation}
  \label{eq:monoexp}
  \boxed{F_k(t) = 1 - \exp\!\left(-\keff^{\text{surf}}\,[E]_t\,t\right)}
\end{equation}
where $\keff^{\text{surf}} \equiv \kcat^{\text{surf}}/K_M^{\text{surf}}$ is the single
measurable surface rate constant. This result is notable: the unknowns $k_1$, $k_{-1}$,
and $\kcat^{\text{surf}}$ individually are inaccessible at clinical $[E]_t$, but
the product $\keff^{\text{surf}}[E]_t$ is directly fit from the progress curve. The
complex accumulates at sub-ppm levels ($\GES/\Gz < 10^{-3}$) and is kinetically invisible.

The solution upper bound is $\keff^{\text{surf}} \leq \kcat^{\text{sol}}/K_m^{\text{sol}}
= 67{,}500\ \Mms$~\cit{2}. Surface immobilisation reduces this by $2$--$100\times$
depending on steric accessibility (Layer~8); the surface value must be measured on the
fabricated electrode.

%=======================================================================
\section{Dimensionless Analysis}
\label{sec:dimensionless}

Three dimensionless numbers characterise the device operating regime and validate the
sequential layer architecture.

\subsection{Kinetic zone parameter $\lambda$}

From Section~\ref{sec:regime}: $\lambda \approx 1$ for all clinical $[E]_t$.
The monoexponential Eq.~\eqref{eq:monoexp} is the appropriate model.
The zone parameter provides an experimental diagnostic: if the measured progress
curve shows $t_{1/2}/t_{1/4}$ significantly below 2.41, the system has left the LG
zone and the full exact solution Eq.~\eqref{eq:theta_exact} must be used.

\subsection{Damk\"{o}hler number: transport does not limit}

The Damk\"{o}hler number compares the surface reaction rate to the enzyme delivery
rate by diffusion~\cit{4}:
\begin{equation}
  \label{eq:Da}
  \Da = \frac{\keff^{\text{surf}}\,[E]_t\,h}{D_E}
\end{equation}
where $h \approx 100\ \mu$m is the diffusion layer height (Layer~3) and
$D_E \approx 4.7\times10^{-7}$~cm$^2$s$^{-1}$ is the MMP-9 diffusion coefficient
(Stokes--Einstein, Layer~3). At the solution upper bound with $[E]_t = 2.17$~nM:
\begin{equation}
  \label{eq:Da_eval}
  \Da^{\text{max}} = \frac{67{,}500 \times 2.17\times10^{-9} \times 10^{-2}}{4.7\times10^{-7}}
  \approx 3.1\times10^{-3} \ll 1
\end{equation}
The electrode surface sees the bulk enzyme concentration to within $<0.4\%$ at all times.
For any plausible surface $\keff^{\text{surf}} < 67{,}500\ \Mms$, $\Da$ is smaller still.
The Layer~3 mass-transport correction is negligible; $[E]_t$ in Eq.~\eqref{eq:monoexp}
is the bulk concentration from Layer~0 throughout.

The Damk\"{o}hler criterion is consistent with the general analysis of Squires
et~al.~\cit{4}: for $\Da \ll 1$ the sensor operates in the reaction-limited regime
where signal is proportional to the surface reaction rate, not the diffusion flux,
and the response is linear in enzyme concentration at short times.

\subsection{Enzyme conservation: $[E]_t$ is constant}

MMP-9 is a true catalyst: it is released intact after each cleavage event
(Eq.~\eqref{eq:mechanism}). The bulk enzyme concentration $[E]_t$ therefore does
not decrease as $\GP$ increases. This is distinct from, for example, irreversible
inhibitors. The constancy of $[E]_t$ is a consequence of the catalytic mechanism,
not an approximation, and holds for all $t$ regardless of $F_k$.

%=======================================================================
\section{Coupled PDE--ODE Numerical Treatment}
\label{sec:coupled}

\subsection{Full coupled system}

For completeness, and to verify the $\Da \ll 1$ approximation, the full coupled
system -- Layer~3 diffusion PDE and Layer~9 surface ODE -- is solved simultaneously:
\begin{align}
  \frac{\partial [E]}{\partial t} &= D_E\,\frac{\partial^2[E]}{\partial z^2},
    \quad 0 < z < h \label{eq:pde} \\
  \frac{d\GES}{dt} &= k_1^*\Gstar - (k_1^* + k_{-1} + \kcat)\GES
    \quad \text{with } k_1^* = k_1[E](z{=}0,t) \label{eq:ode_coupled}
\end{align}
with boundary conditions:
\begin{equation}
  \left.D_E\,\frac{\partial[E]}{\partial z}\right|_{z=0} = \kcat\,\GES(t)
  \quad\text{(flux = cleavage consumption)}, \qquad
  [E](h,t) = [E]_{\text{bulk}}
  \label{eq:bc}
\end{equation}
Solved numerically: Crank--Nicolson finite difference for the PDE, RK4 for the surface
ODE, coupled at each timestep through the $z=0$ boundary condition.

\subsection{Verification of the decoupled limit}

For $\Da \ll 1$, the PDE solution gives $[E](0,t) \approx [E]_{\text{bulk}}$ to
within $\Da/(1+\Da) < 0.4\%$. The surface ODE then has a constant-coefficient
enzyme input and reduces exactly to Eq.~\eqref{eq:monoexp}. The coupled simulation
(\texttt{surface\_kinetics.py}) confirms this: the full PDE--ODE $F_k(t)$ is
indistinguishable from the analytic monoexponential at all clinical $[E]_t$ values.
The decoupled monoexponential is therefore used for all subsequent analysis.

%=======================================================================
\section{$F_k(t)$ Output and Optimal Readout Time}
\label{sec:Fk}

\subsection{Progress curves}

From Eq.~\eqref{eq:monoexp}, the fractional cleavage at incubation time $T$:
\begin{equation}
  \label{eq:Fk}
  F_k(T) = 1 - \exp\!\left(-\keff^{\text{surf}}\,[E]_t\,T\right)
\end{equation}
In the linear regime ($\keff^{\text{surf}}[E]_t T \ll 1$):
\begin{equation}
  \label{eq:Fk_linear}
  F_k(T) \approx \keff^{\text{surf}}\,[E]_t\,T
\end{equation}
The ratio of chronic to healing fractional cleavage is:
\begin{equation}
  \label{eq:ratio_linear}
  \frac{F_k^{\text{chronic}}}{F_k^{\text{healing}}} \approx
  \frac{[E]_{\text{chronic}}}{[E]_{\text{healing}}} = \frac{30\ \text{nM}}{6\ \text{nM}} = 5\times
\end{equation}
independent of $\keff^{\text{surf}}$ and $T$. The 5$\times$ discrimination ratio
is therefore robust across the entire plausible range of surface rate constants.

\subsection{Optimal readout time $T^*$}

In the linear regime, $F_k \propto [E]_t T$, so the ratio $F_k^c/F_k^h$ is
always $5\times$ regardless of $T$. The optimal incubation time is then as long as
possible before signal saturation ($F_k \to 1$).

In the exponential (nonlinear) regime, the \emph{absolute} signal difference
$\Delta F_k = F_k^{\text{chronic}} - F_k^{\text{healing}}$ varies with $T$.
Maximising $\Delta F_k$ gives the optimal readout time. Setting
$d(\Delta F_k)/dT = 0$:
\begin{align}
  \keff^{\text{surf}}\,[E]_c\,\exp\!\left(-\keff^{\text{surf}}[E]_c T\right)
  &= \keff^{\text{surf}}\,[E]_h\,\exp\!\left(-\keff^{\text{surf}}[E]_h T\right) \nonumber\\
  \Rightarrow\quad
  T^* &= \frac{\ln([E]_c/[E]_h)}{\keff^{\text{surf}}\,([E]_c - [E]_h)}
  \label{eq:Tstar}
\end{align}
With $[E]_c = 30$~nM, $[E]_h = 6$~nM, and ratio $r = 5$:
\begin{equation}
  \label{eq:Tstar_eval}
  T^* = \frac{\ln 5}{\keff^{\text{surf}} \times 24\ \text{nM}}
  = \frac{1.609}{\keff^{\text{surf}} \times 2.4\times10^{-8}\ \text{M}}
\end{equation}
\begin{center}
\begin{tabular}{lll}
\toprule
$\keff^{\text{surf}}$ (M$^{-1}$s$^{-1}$) & $T^*$ (s) & $T^*$ (min) \\
\midrule
$2{,}000$    & $33{,}500$ & 558 \\
$10{,}000$   & $6{,}700$  & 112 \\
$30{,}000$   & $2{,}230$  &  37 \\
$67{,}500$ (sol.\ UB) & $990$ & 16 \\
\bottomrule
\end{tabular}
\end{center}
The target incubation window of 30--60~min corresponds to
$\keff^{\text{surf}} \approx 13{,}000$--$27{,}000\ \Mms$. If the measured
$\keff^{\text{surf}}$ falls significantly below this, the incubation time must be
extended accordingly.

%=======================================================================
\section{Temperature Dependence}
\label{sec:temperature}

\subsection{Arrhenius correction}

The surface rate constant follows Arrhenius temperature dependence:
\begin{equation}
  \label{eq:arrhenius}
  \keff^{\text{surf}}(T_{\text{amb}}) = \keff^{\text{surf}}(T_{\text{ref}})\,
  \exp\!\left[-\frac{E_a}{R}\left(\frac{1}{T_{\text{amb}}} - \frac{1}{T_{\text{ref}}}\right)\right]
\end{equation}
where $E_a$ is the activation energy and $T_{\text{ref}} = 310$~K (37$^\circ$C).
The relative sensitivity to ambient temperature is:
\begin{equation}
  \label{eq:arrhenius_sensitivity}
  \frac{1}{\keff}\frac{d\keff}{dT_{\text{amb}}} = \frac{E_a}{R\,T_{\text{amb}}^2}
\end{equation}
\textit{The activation energy $E_a$ for MMP-9 surface cleavage of SGKIPRTLTA is not
confirmed from any paper in \texttt{papers/layer\_09/}.} Using a typical
metalloproteinase value $E_a \approx 45$~kJ~mol$^{-1}$ (unconfirmed placeholder):
\begin{equation}
  \left.\frac{d\ln\keff}{dT}\right|_{310\,\text{K}}
  \approx \frac{45{,}000}{8.314 \times 310^2} \approx 5.6\%\ \text{per}\ {}^\circ\text{C}
  \label{eq:temp_sensitivity}
\end{equation}

\subsection{Propagation to $F_k$}

Let $x = \keff^{\text{surf}}[E]_t T$. The temperature sensitivity of $F_k$ is:
\begin{equation}
  \label{eq:dFdT}
  \frac{\partial F_k}{\partial T_{\text{amb}}} = (1-F_k)\,x\,\frac{E_a}{R\,T_{\text{amb}}^2}
  = (1-F_k)\,\left|\ln(1-F_k)\right|\,\frac{E_a}{R\,T_{\text{amb}}^2}
\end{equation}
At $F_k = 0.1$ (early signal, $x = 0.105$) and $\pm 1^\circ$C uncertainty:
$\delta F_k \approx 0.1 \times 0.105 \times 0.056 \approx 6\times10^{-4}$,
a 0.6\% absolute error on $F_k$. At $F_k = 0.5$ ($x = 0.693$):
$\delta F_k \approx 0.5 \times 0.693 \times 0.056 \approx 1.9\%$.

Temperature control or measurement is required if absolute $F_k$ accuracy
$< 2\%$ is needed. The ratiometric signal $\Delta R/R_0$ is similarly affected
since $F_k$ enters linearly (Layer~13). The discrimination \emph{ratio}
$F_k^c/F_k^h$ is less sensitive: at short times both $F_k^c$ and $F_k^h$ scale
with the same $\keff(T)$, so the ratio $[E]_c/[E]_h = 5\times$ is temperature-invariant.

%=======================================================================
\section{Multi-Protease Competition}
\label{sec:competition}

\subsection{General competition model}

Wound fluid contains multiple proteases that may cleave SGKIPRTLTA. In the first-order
(linear) regime, contributions add independently:
\begin{equation}
  \label{eq:competition}
  \frac{dF_k}{dt} = k_{\text{obs,total}} = \sum_j k_{\text{eff},j}^{\text{surf}}\,[E_j]_t
\end{equation}
The fractional contribution of interferent $j$ to the observed cleavage signal is:
\begin{equation}
  \label{eq:interferent}
  \xi_j = \frac{k_{\text{eff},j}^{\text{surf}}\,[E_j]_t}{k_{\text{obs,total}}}
\end{equation}
MMP-9 selectivity requires $\xi_{\text{MMP-9}} \approx 1$, i.e.\
$k_{\text{eff,MMP-9}}^{\text{surf}}[E_{\text{MMP-9}}] \gg \sum_{j \neq \text{MMP-9}}
k_{\text{eff},j}^{\text{surf}}[E_j]$.

\subsection{Known interferents and status}

\begin{tabular}{llll}
\toprule
Interferent & $[E_j]$ source & $k_{\text{eff},j}^{\text{surf}}$ status & Risk \\
\midrule
HNE    & 30~nM chronic (Layer~0, Trengove~1999) & Unconfirmed & Moderate \\
MMP-2  & Wound fluid present, conc.\ unconfirmed & Unconfirmed & Low--moderate \\
MMP-8  & Wound fluid present, conc.\ unconfirmed & Unconfirmed & Low \\
\bottomrule
\end{tabular}

\medskip
\textit{No paper in \texttt{papers/layer\_09/} provides $k_{\text{eff}}^{\text{surf}}$ for any
interferent protease on SGKIPRTLTA. The competition fractions $\xi_j$ cannot be
computed until these are measured.} HNE is the highest-priority interferent given its
confirmed elevated concentration in chronic VLU and its Ser-protease activity
overlapping with some MMP substrates.

%=======================================================================
\section{Uncertainty Propagation}
\label{sec:uncertainty}

\subsection{Analytic uncertainty in $F_k$}

Let $x = \keff^{\text{surf}}[E]_t T$, so $F_k = 1 - e^{-x}$. The partial
derivatives are:
\begin{equation}
  \label{eq:partials}
  \frac{\partial F_k}{\partial \keff} = \frac{x}{\keff}\,e^{-x}, \qquad
  \frac{\partial F_k}{\partial [E]_t} = \frac{x}{[E]_t}\,e^{-x}, \qquad
  \frac{\partial F_k}{\partial T} = \frac{x}{T}\,e^{-x}
\end{equation}
By standard propagation:
\begin{equation}
  \label{eq:sigma_Fk}
  \sigma_{F_k} = e^{-x}\,x\;\sqrt{
    \left(\frac{\sigma_{\keff}}{\keff}\right)^2
    +\left(\frac{\sigma_{[E]}}{[E]_t}\right)^2
    +\left(\frac{\sigma_T}{T}\right)^2}
  = (1-F_k)\,\left|\ln(1-F_k)\right|\,\sigma_{\text{rel,total}}
\end{equation}

\subsection{Dominant uncertainty sources}

\begin{tabular}{lll}
\toprule
Source & Relative uncertainty $\sigma/\mu$ & Status \\
\midrule
$\keff^{\text{surf}}$ & Unknown; plausible range $2{,}000$--$67{,}500\ \Mms$ (factor $34\times$) & Dominant until measured \\
$[E]_t$ & $\sim15$--$20\%$ (Layer~0 ODE; healing estimate approximate) & Secondary \\
$T$ (incubation) & $\sim1$--$2\%$ (timer precision) & Negligible \\
$T_{\text{amb}}$ & $\sim0.6\%$ per ${}^\circ$C (Eq.~\ref{eq:temp_sensitivity}) & Tertiary \\
\bottomrule
\end{tabular}

\medskip
The dominant uncertainty is $\keff^{\text{surf}}$, which spans more than an order of
magnitude before electrode calibration. All $F_k$ predictions are therefore presented
as parametric sweeps over the plausible $\keff^{\text{surf}}$ range. Once the
fabricated electrode is measured (Layer~15 calibration), $\keff^{\text{surf}}$ is
pinned to $\pm 5$--$10\%$ and $\sigma_{F_k}$ is determined by $\sigma_{[E]}$ alone.

%=======================================================================
\section{Parameter Table and Layer Outputs}
\label{sec:params}

\rowcolors{2}{gray!8}{white}
\begin{tabularx}{\linewidth}{lXlll}
\toprule
Parameter & Description & Value & Source & Confirmed \\
\midrule
$\kcat^{\text{sol}}$ & Solution-phase catalytic rate for SGKIPRTLTA & 86.4~s$^{-1}$ & Kridel~2001~\cit{2} Table~II & Yes \\
$K_m^{\text{sol}}$ & Solution-phase Michaelis constant & 1.28~mM & Kridel~2001~\cit{2} Table~II & Yes \\
$(\kcat/K_m)^{\text{sol}}$ & Solution efficiency (UB for surface) & $67{,}500\ \Mms$ & Kridel~2001~\cit{2} Table~II & Yes \\
$[E]_t^{\text{chronic}}$ & Free MMP-9 at 30-min read, chronic & 30~nM & Layer~0 ODE & Yes \\
$[E]_t^{\text{healing}}$ & Free MMP-9 at 30-min read, healing & 6~nM & Layer~0 ODE & Yes \\
$\Gz$ & Initial surface probe density & $\sim$5~pmol~cm$^{-2}$ & Layer~7 (device-planned) & No \\
$\keff^{\text{surf}}$ & Surface rate constant ($\kcat/K_m^{\text{sol}}$, upper bound) & To be measured & Layer~15 calibration & No \\
$K_{m,\text{surf}}^\Gamma$ & Surface MM constant (mol~cm$^{-2}$) & $6.4\times10^{-13}$~mol~cm$^{-2}$ (estimate, Eq.~\ref{eq:Kmsurf_val}) & Eq.~\ref{eq:Kmsurf}; Layer~8 $\delta_{\text{eff}}$ unconfirmed & No \\
$\rho$ & Probe-saturation ratio $\Gz/K_{m,\text{surf}}^\Gamma$ & 0.3--3.1 (Table~\ref{tab:rho}) & Layer~8 steric factor & No \\
$E_a$ & Activation energy for surface cleavage & $\sim$45~kJ~mol$^{-1}$ placeholder & Not in \texttt{papers/layer\_09/} & No \\
$A$ & Electrode area & $\sim$0.002~cm$^2$ & Device-planned & No \\
$k_{\text{eff},j}^{\text{surf}}$ & Surface rate constant for interferents & Unknown & Not in \texttt{papers/layer\_09/} & No \\
\bottomrule
\end{tabularx}

\medskip
\noindent\textbf{Layer outputs.}
\begin{itemize}[noitemsep]
  \item $F_k(T)$: fractional cleavage as function of incubation time $\to$ Layer~13 (ratiometric signal $\Delta R/R_0$)
  \item $k_{\text{obs}} = \keff^{\text{surf}}[E]_t$: pseudo-first-order rate $\to$ Layer~14 (noise budget)
  \item $T^*$: optimal incubation time (Eq.~\eqref{eq:Tstar}) $\to$ device protocol
  \item $\keff^{\text{surf}}$ (post-calibration) $\to$ Layer~1 (dilution factor), Layer~15 (calibration curve)
\end{itemize}

%=======================================================================
\section{Limitations}
\label{sec:limitations}

\begin{enumerate}[leftmargin=*]

  \item \textbf{$\keff^{\text{surf}}$ is unknown; all $F_k$ predictions are parametric.}
    The surface rate constant $\keff^{\text{surf}}$ must be measured on the fabricated
    electrode before any absolute signal or LOD can be specified. The solution upper bound
    $\kcat^{\text{sol}}/K_m^{\text{sol}} = 67{,}500\ \Mms$~\cit{2} is the starting
    assumption; surface immobilisation typically reduces this by $2$--$100\times$
    (Layer~8 steric correction, Adjémian~2010~\cit{3}). Until calibration, the plausible
    range spans $2{,}000$--$67{,}500\ \Mms$ and $T^*$ spans 16~min to $>9$~h.

  \item \textbf{$K_M^{\text{surf}}$ and the LG zone assumption require experimental
    verification.} The $\lambda \approx 1$ conclusion holds for any $k_1 \leq 10^8\ \Mms$
    and $[E]_t \leq 2.17$~nM. However, if the effective surface $K_M^{\text{surf}}$ is
    pushed above $\sim 100$~nM by steric crowding (putting $[E]_t/K_M^{\text{surf}}$
    above 0.02), the system could enter a regime where the exact biexponential
    Eq.~\eqref{eq:theta_exact} is needed for accurate fitting. The diagnostic ratio
    $t_{1/2}/t_{1/4} \to 2.41$ (LG) vs.\ $\to 1.75$ (FS) should be measured~\cit{1}.

  \item \textbf{Temperature dependence is unconfirmed.}
    The Arrhenius activation energy $E_a$ for MMP-9 cleavage of SGKIPRTLTA on a
    surface is not reported in any paper in \texttt{papers/layer\_09/}. The 5.6\%~per
    ${}^\circ$C figure is a placeholder based on typical metalloproteinase values. A
    temperature-controlled incubation or in-situ temperature measurement is required
    for absolute $F_k$ accuracy.

  \item \textbf{Interferent cross-reactivity is unquantified.}
    The surface rate constants $k_{\text{eff},j}^{\text{surf}}$ for HNE, MMP-2, MMP-8, and
    other wound-fluid proteases on SGKIPRTLTA are not reported in any paper in
    \texttt{papers/layer\_09/}. The selectivity fractions $\xi_j$ cannot be computed.
    HNE is the priority interferent given its confirmed elevated chronic wound
    concentration (30~nM, Layer~0~\cit{2}) and potential serine-protease activity on
    the substrate.

  \item \textbf{Solution-to-surface $K_m$ conversion requires $\delta_{\text{eff}}$
    from Layer~8.}
    The effective surface Michaelis constant $K_{m,\text{surf}}^\Gamma$ is related to the
    solution value by $K_{m,\text{surf}}^\Gamma = K_m^{\text{sol}} \times 10^{-3} \times
    \delta_{\text{eff}}$, where $\delta_{\text{eff}}$ (Layer~8, unconfirmed) is the effective
    thickness of the accessible probe layer. Until $\delta_{\text{eff}}$ is established from
    Layer~8, $K_{m,\text{surf}}^\Gamma$ spans approximately $6\times10^{-14}$--$6\times10^{-12}$~mol~cm$^{-2}$
    ($\delta_{\text{eff}}$ range 0.5--50~nm).

  \item \textbf{The probe-saturation ratio $\rho$ is unconfirmed; the
    progress curve shape is therefore uncertain.}
    Table~\ref{tab:rho} shows $\rho \in [0.3, 3.1]$ across the plausible
    steric correction range. Over this range, the progress curve transitions
    from near-monoexponential ($\rho \leq 0.3$) to significantly non-monoexponential
    ($\rho \geq 1$). The implicit integrated equation
    (Eq.~\eqref{eq:imp_mm}) must be fit to measured progress curves to determine
    $\rho$ empirically; fitting with the monoexponential alone will overestimate
    $F_k$ at intermediate times if $\rho > 0.3$. This is the dominant source of
    \emph{systematic} modelling error until $\Gz$, $\delta_{\text{eff}}$, and
    $\kcat$ are measured on the fabricated electrode.

\end{enumerate}

%=======================================================================
\section*{References}

\noindent All papers in \texttt{electrochemical\_device/papers/layer\_09/}.

\begin{enumerate}[label={[\arabic*]},noitemsep,leftmargin=*]

\item Anne A, Demaille C (2012). Kinetics of enzyme action on surface-attached substrates:
  a practical guide to progress curve analysis in any kinetic situation.
  \textit{Langmuir} 28:14665--14671.
  \textbf{anne\_demaille2012.pdf} -- ODE system (Eq.~1), exact biexponential solution
  (Eqs.~4--5), $\lambda$ parameter (Eq.~6 region), universal progress curve (Eq.~7),
  FS limit (Eq.~8), zone diagram (Fig.~4): all confirmed from text.

\item Kridel SJ, Chen E, Howard EW, Bhatt L, Strongin AY (2001). Substrate hydrolysis
  by matrix metalloproteinase-9.
  \textit{J Biol Chem} 276:20572--20578.
  \textbf{kridel2001.pdf} -- $\kcat = 86.4$~s$^{-1}$, $K_m = 1.28$~mM,
  $\kcat/K_m = 67{,}500\ \Mms$ for SGKIPRTLTA (peptide A11), Table~II confirmed.

\item Adjémian J, Anne A, Cauet G, Demaille C (2010). Cleavage-sensing peptide-enzyme
  multilayer gold surfaces for the parallel, selective, and sensitive assay of trypsin
  and plasmin activities by electrochemical methods.
  \textit{Langmuir} 26:10347--10356.
  \textbf{adjemian2010.pdf} -- experimental surface cleavage yields and $\Gz$ context.

\item Squires TM, Messinger RJ, Manalis SR (2008). Making it stick: convection,
  reaction and diffusion in surface-based biosensors.
  \textit{Nat Biotechnol} 26:417--426.
  \textbf{squires2008.pdf} -- Damk\"{o}hler number framework for reaction-limited
  vs.\ transport-limited biosensor regimes; surface MM ODE formulation.

\item Pheeney CG, Barton JK (2012). DNA electrochemistry with tethered methylene blue.
  \textit{Langmuir} 28:7063--7070. DOI:~10.1021/la300566x.
  \textbf{(in papers/layer\_10/)} -- MB surface coverage
  $\Gz = 1.3$--$3.2$~pmol~cm$^{-2}$ used for $\rho$ calculation in
  Table~\ref{tab:rho}.

\end{enumerate}

\end{document}
